% Placement: methods_rewrite.tex, end of \subsubsection*{Proportional--integral feedback},
% after the paragraph ending "...would describe a controller other than the one that ran." (L323).
% Uses the existing rem environment (\newtheorem{rem}{Remark}, main.tex L25).

The structure of the controller nonetheless rests on a standard result for the nominal, trial-averaged plant.

\begin{rem}[Low-gain integral tracking through a monotone actuator]
\label{rem:lowgain}
Write the nominal plant from drive to suppression $s = -y$ as a linear time-invariant system $G$ driven through a static nonlinearity $\phi$ of the integrator state $u$. With $K_p = 0$ and $u = K_i\bar e$, the controller of Eq.~\eqref{eq:pi} is the loop below with $k = K_i$, provided $\phi$ absorbs the feedforward offset $\bar u$, the clamp of the integral term to $[-q, q]$, the clipping of the command to $[0, u_{\max}]$ and the relation between command and effective drive. Suppose that (i) $G$ is exponentially stable, possibly with an input delay; (ii) its steady-state gain is positive, $G(0) > 0$; (iii) $\phi$ is nondecreasing and globally Lipschitz with constant $\lambda$ (we ignore the 0.01\,V quantization of the command, a staircase that is not Lipschitz); and (iv) the reference is feasible, that is, $|y_{\mathrm{ref}}|/G(0)$ lies in the range of $\phi$, which with the clamps above is the drive delivered by commands between $\max(0, \bar u - q)$ and $\min(u_{\max}, \bar u + q)$; if the unstimulated output carries a constant offset $y_0$, $|y_{\mathrm{ref}}|$ is replaced by $|y_{\mathrm{ref}}| + y_0$. Then there is a $K > 0$, depending only on $G$, such that for every integral gain $0 < k < K/\lambda$ the integral loop $\dot{u} = k\,e$, with $e = |y_{\mathrm{ref}}| - s = y - y_{\mathrm{ref}}$ as in Eq.~\eqref{eq:error}, tracks the reference: $\phi(u(t)) \to |y_{\mathrm{ref}}|/G(0)$, $u(t)$ converges, the error is square-integrable, and $e(t) \to 0$ under a regularity condition that is mild for a finite-dimensional plant with input delay \citep[Theorem~4.1(a),(c),(d),(f) and Remark~4.2(b)]{logemann2000absolute}. A gain below which tracking is guaranteed, that is, a lower bound on $K/\lambda$, can be computed from open-loop step- or frequency-response data \citep{logemann1999integral}; analogous results hold in discrete time \citep{fliegner2001discrete} and, for proportional--integral control, for nonlinear plants whose steady-state input--output map is an increasing bijection with slope bounded away from zero, a condition that a saturating actuator violates \citep{desoer1985tracking}.
\end{rem}

The trial-averaged impulse-response models (Fig.~\ref{fig:figure2}), recorded in separate sessions and in two of the four controller mice, are stable, satisfying (i). They do not by themselves fix the sign of $G(0)$, the integral of the impulse response over all time, because their slowest mode exceeded the 0.5\,s fit window in three of four sessions. That the response keeps its suppressive sign over a trial, which is what (ii) requires on our time scale, is shown instead by the open-loop trials, in which a constant command held activity below baseline for the whole trial (Fig.~\ref{fig:figure3}B). Evoked suppression grew monotonically with laser power (Fig.~\ref{fig:figure2}B); although it is a transient measure, it makes the steady-state map monotone through the structure of the nominal model, a static actuator map feeding linear dynamics, given (ii). The clipping, the clamp on the integral term and this monotone dose--response make $\phi$ nondecreasing, satisfying (iii), and (iv) requires the reference, plus any standing offset of the unstimulated activity, to lie within the suppression attainable with the commands the clamps allow. The controller that ran departs from the remark in four further ways: it had a proportional term, its integral was a 3\,s trailing window rather than an unbounded sum (Eq.~\eqref{eq:integral}), it ran in discrete time at 35\,Hz, and each trial lasted 3\,s (4\,s with the sinusoidal reference) with the integrator reset between trials, whereas the remark is asymptotic and concerns the nominal plant, not one that varies with brain state. The remark therefore motivates the controller's structure rather than certifying the loop that ran. That loop's stability rests on the observation that no session showed sustained oscillation; with a stable plant and a bounded command the output cannot diverge, so oscillation is the instability the data could have shown.
